\chapter{INTRODUCTION}
\pagestyle{fancy}
\pagenumbering{arabic}

\section{Motivation}
In the modern data economy, organizations process millions of events per second through streaming pipelines, where milliseconds matter and decisions propagate in real time. Unlike batch processing---where quality checks can be applied retrospectively after data has been collected---streaming environments demand instantaneous validation as each record arrives. This fundamental difference in timing creates a unique quality challenge: by the time a traditional batch validator detects a completeness violation, the downstream fraud-detection model has already consumed the dirty record, the customer has been incorrectly charged, and the financial report has been updated with corrupted figures. The contamination cascade in streaming pipelines is thus not merely a data hygiene problem but a systemic risk that amplifies errors at the speed of computation.

The financial stakes of this problem are substantial. IBM research estimates that poor data quality costs organizations an average of \textbf{USD 12.9 million annually} \cite{ibm_data_quality_cost}. The Databricks Lakehouse platform documents similar figures, noting that data quality failures directly erode the value proposition of analytical and ML investments \cite{Armbrust2021Lakehouse}. IBM further articulates that the "garbage in, garbage out" principle is especially acute for machine learning: models trained on contaminated data yield predictions that propagate errors through every downstream decision, from automated approvals to revenue forecasting \cite{ibm_dq2024}. In high-stakes domains, the consequences extend beyond cost. Ververica reports that delayed compliance detection in financial services alone accounts for \textbf{\$4.7 billion in regulatory fines annually} \cite{ververica2024}, underscoring that streaming data quality is both a technical and a regulatory imperative.

The operational impact of bad data manifests across the entire data ecosystem. Confluent identifies eight distinct categories of bad data in streaming pipelines, manifesting as inaccurate dashboards, failed compliance audits, customer churn, and wasted compute on poisoned ML models \cite{confluent_2025}. A Vimeo case study documented a one-day lag before insights reached analytics teams, directly limiting real-time campaign pivots \cite{confluent_2025}. Grab Engineering reports that "missing a few insert booking records in peak hours can generate wrong downstream results leading to miscalculation in revenue" \cite{grab_streaming_dq}---a concrete instance where data quality failures translate immediately into business losses. Conversely, Siemens Healthineers processes 8 million messages daily via Confluent, catching manufacturing defects in real time and enabling reliable patient diagnostics \cite{confluent_2025}---evidence that streaming quality validation creates operational and clinical value when implemented correctly.

The contamination cascade is particularly acute in streaming environments. Confluent warns that once bad data enters a stream, "it spreads to all consumers, contaminating every dependent dataset" \cite{confluent_2025}, and remediation requires "surgical removal + reprocessing + downstream recomputation" across the entire downstream ecosystem. Traditional batch-style validation compounds this problem: by the time a nightly quality check flags an error, damage has already propagated to payment processing, fraud detection, and financial reporting \cite{confluent_2025}. The Confluent 2025 Data Streaming Report estimates that only \textbf{~1\% of companies} have reached the highest streaming maturity level, where data streams are treated as managed products with strategic quality guarantees \cite{confluent_2025_streaming_report}. This statistic highlights both the scale of the current gap and the opportunity for frameworks that address streaming data quality systematically.

These challenges motivate our investigation using the NYC Yellow Taxi dataset from the NYC Taxi and Limousine Commission (TLC) as a representative urban mobility use case. Empirical analysis of 72 million records spanning 26 months reveals five categories of data quality problems whose interactions create compounding risks in production pipelines. First, \textbf{completeness violations} occur when required fields are missing or null, affecting approximately 9.1\% of records in our dataset. Second, \textbf{validity violations} arise when individual field values violate physical constraints---such as negative fares, zero distances, or impossible speeds---impacting roughly 3.4\% of records. Third, \textbf{multivariate anomalies} represent the most insidious category: combinations of individually reasonable feature values are collectively suspicious and require understanding inter-feature correlations. These patterns are undetectable by single-feature rules and constitute approximately 2.9\% of anomalous records. Fourth, \textbf{uniqueness violations} emerge from duplicate records appearing due to data reprocessing. Fifth, \textbf{concept drift} causes underlying data distributions to shift over time, rendering previously calibrated thresholds and trained models stale. These five problem categories interact: a drift in passenger pickup zones may alter the expected distribution of fares, causing validity thresholds calibrated for one zone distribution to generate false positives in another. This interaction is precisely why existing single-strategy frameworks fail---each category demands distinct detection and remediation approaches, and no existing framework addresses all five jointly.

\section{Research Challenges}
Despite significant research progress in data quality management, three critical challenges remain unresolved in existing streaming DQ frameworks.

The first challenge is \textbf{context collapse in heterogeneous data}. Traditional frameworks apply a single global threshold to all records regardless of context, producing catastrophically high false positive rates on heterogeneous data. Consider the NYC Taxi dataset, which exhibits extreme heterogeneity across multiple dimensions: trip type (airport versus street-hail), time of day (rush hour versus late night), day of week (weekday versus weekend), and geographic zone (Manhattan versus outer boroughs). A global threshold of fare greater than \$45, calibrated for standard Manhattan trips, flags legitimate Newark airport fares---with a mean of \$68.44 and a standard deviation that places the majority of airport trips above the threshold---as anomalies at rates up to 68.9\%. Conversely, this threshold remains too permissive for short-distance trips, where fraudulent transactions at \$40-\$45 may go undetected. The DTD framework formally proves that no single fixed threshold is universally optimal across heterogeneous contexts \cite{dtd2026}; yet existing streaming DQ systems universally apply global thresholds. A systematic literature review on context in data quality management \cite{contextdq2022} confirms that existing tools (Confluent data contracts, Grab Coban, Datadog anomaly detection) focus primarily on schema validation and rule-based checks without incorporating multi-dimensional contextual signals (geographic location, time-of-day, day-of-week, zone type) into real-time quality scoring. Memory-based streaming anomaly detectors such as MemStream \cite{bhatia2022memstream} and ADWIN-U \cite{adwin2025} maintain historical data distributions but do not stratify quality expectations by geographic zone or temporal period, despite evidence that spatial and temporal context significantly improves anomaly detection in urban mobility data \cite{bestad2025}.

The second challenge is \textbf{sequential pipeline architectures that underutilize heterogeneous validation complexity}. Existing streaming DQ frameworks process records through validation stages sequentially: Schema Validation, followed by Business Rules, then ML scoring, and finally Drift Detection. This linear architecture compounds end-to-end latency and wastes computational resources on records that earlier stages could flag as violations. The fork-join pattern offers a natural solution: records that fail early structural checks (schema, nullity, range validity) can bypass expensive ML inference entirely, while records passing early checks proceed to behavioral analysis. Approximately 13.5\% of input records in a well-designed linear pipeline fail at the Schema or Business Rules stage---a proportion that motivates the fork-join design but is not an empirical constant; it varies by dataset and validation rules. Existing frameworks do not implement this branching: industry tools rely on rule-based validation while academic streaming anomaly detectors rely on ML-based scoring, and no existing work synthesizes both branches within a unified streaming pipeline \cite{contextdq2022}.

The third challenge is \textbf{single-strategy drift adaptation that fails to differentiate drift severity}. When concept drift is detected, existing frameworks apply a uniform response: full model retraining regardless of the magnitude or type of the detected shift. This one-size-fits-all approach wastes compute on minor distributional fluctuations that lightweight threshold updates could resolve, while potentially under-responding to severe distribution shifts that require comprehensive model rebuilding. The DTD framework \cite{dtd2026} formally proves that no single fixed adaptation strategy is universally optimal across drift types. NYC Taxi data exhibits multiple drift types: gradual seasonal shifts (summer tourism patterns), sudden event-driven spikes (weather events, holidays), and recurring cyclical patterns (weekly ridership cycles) that demand differentiated responses. A minor gradual shift in zone-level fare distributions may be handled by updating per-zone thresholds incrementally, while a sudden spike caused by a fare structure change may require full model warmup on fresh data. Existing streaming DQ systems lack automated strategy selection for quality signals \cite{contextdq2022}, and real-time meta-aggregation for drift monitoring across quality dimensions remains underexplored. Additionally, adaptive per-zone quality thresholds are not well studied for streaming data. ReAD \cite{read2020} demonstrates that dynamic, region-specific anomaly thresholds outperform fixed global thresholds in spatial anomaly detection, but no existing work applies this insight to streaming data quality evaluation for urban mobility applications.

\section{Objective}
The primary objective of this thesis is to design and implement \textbf{CA-DQStream} (Context-Aware Data Quality for Streaming), a comprehensive framework that addresses the five data quality problem categories in streaming environments through context-aware, adaptive, and parallel architectures. Specifically, this thesis pursues five sub-objectives:

\begin{enumerate}
    \item Design and implement a five-layer streaming DQ architecture spanning Data Source, Kafka ingestion, Apache Flink processing, data storage, and monitoring layers.
    \item Develop a four-stage processing pipeline combining Schema Validation, Business Rules, Online Autoencoder with Memory Module anomaly detection, and ADWIN-based drift detection.
    \item Develop a self-monitoring mechanism that uses ADWIN drift detection on quality meta-streams to detect when per-zone baselines become stale.
    \item Evaluate CA-DQStream on the NYC Yellow Taxi dataset, demonstrating improved false positive rate control across heterogeneous zones and effective multi-type drift detection.
    \item Implement an MLOps pipeline integrating Online AE warmup, Flink Broadcast State deployment for zone-aware thresholds, and IEC-triggered memory module updates.
\end{enumerate}

CA-DQStream is specifically suited for urban mobility streaming data where contextual variation (rush hour versus late night, residential versus commercial zones, weekday versus weekend) fundamentally changes what constitutes high-quality data. By introducing zone-aware and periodicity-aware quality profiles, CA-DQStream addresses the intersection of context-awareness, quality-aware scoring, and adaptive thresholds identified in the systematic literature review on context-aware DQ \cite{contextdq2022}.

\section{Contribution Highlights}
CA-DQStream makes four core scientific contributions that address the identified research challenges:

\begin{itemize}
    \item \textbf{4D Context-Aware Thresholding:} \textit{Problem:} Global anomaly thresholds produce excessive false positives in heterogeneous urban mobility data, where legitimate high-fare airport trips are misclassified as anomalies. \textit{Insight:} Contextual variation across geographic zones, temporal windows, trip types, and day categories is the primary driver of threshold sensitivity; a matrix of context-specific thresholds can capture this variation. \textit{Method:} A four-dimensional threshold matrix (trip\_type $\times$ time\_window $\times$ day\_type $\times$ neighborhood\_id) with 588 cells, computed automatically as $\mu_{\text{cell}} + 2\sigma_{\text{cell}}$ from training data with intelligent fallback for sparse cells. The 588 cells arise from a grid decomposition of the NYC Taxi dataset's natural heterogeneity: 3 trip types (airport, street-hail, shared), 7 time windows (hour-of-day blocks), 4 day types (weekday/weekend $\times$ rush-hour/off-peak), and 7 geographic zones (manhattan zones 1-6 plus outer boroughs). This 4D decomposition captures the empirically observed variance structure where zone explains 34\% of fare variance and time-of-day explains an additional 18\%. \textit{Validation:} This approach directly addresses the DTD theorem's conclusion that no single fixed threshold is universally optimal \cite{dtd2026}. ReAD \cite{read2020} independently validates that dynamic, region-specific anomaly thresholds outperform fixed global thresholds in spatial anomaly detection.

    \item \textbf{Rendezvous Pipeline Architecture:} \textit{Problem:} Sequential validation pipelines compound end-to-end latency and waste ML inference resources on records that fail early structural validation. \textit{Insight:} A fork-join architecture can route records through the appropriate validation complexity level based on early-stage pass/fail signals, enabling cheap records to bypass expensive ML scoring. \textit{Method:} A fork-join pipeline with two parallel branches: a Canary branch (schema validation, null checks, range validity) and a Complex branch (ML-based autoencoder anomaly scoring, multivariate correlation checks). Both branches converge at a MetaAggregator rendezvous point---a synchronization barrier that collects validation outputs from both branches, resolves conflicts, and emits a unified quality verdict per record. Records failing Canary checks exit at the rendezvous point without entering the Complex branch, reducing mean end-to-end latency. Records passing Canary checks proceed to the Complex branch in parallel, enabling throughput improvements through concurrent execution. \textit{Validation:} This architecture resolves the dual-branch synthesis gap identified in the literature \cite{contextdq2022} by combining rule-based structural quality checks with ML-based behavioral anomaly detection within a unified streaming pipeline, informed by Grab's FlinkSQL approach for real-time quality monitoring \cite{grab_coban2025}.

    \item \textbf{Intelligent Evolution Controller (IEC) with Four Adaptation Strategies:} \textit{Problem:} Concept drift in streaming data causes trained quality models to become stale, but uniform retraining responses waste compute on minor shifts and under-respond to severe ones. \textit{Insight:} Drift events vary in severity and type, and automated strategy selection based on drift characteristics can match adaptation intensity to drift magnitude. \textit{Method:} A multi-strategy self-adaptation mechanism that monitors quality meta-streams via ADWIN and selects from four complementary strategies: Continuous Evolution (incremental threshold updates for gradual drift), Switching Scheme (model checkpoint switching for sudden drift), METER Adaptation (hypernetwork-based centroid prediction for cyclical drift), and Spatial Tracking (per-zone threshold recalibration for zone-specific drift). The meta-aggregation layer monitors drift in four quality score distributions (completeness, timeliness, consistency, validity) across reference and current windows, detecting when the ADWIN window divergence exceeds a configurable threshold and triggering the appropriate adaptation strategy. This hierarchy enables 93\% of drift events to be resolved through lightweight incremental updates rather than expensive full retraining. \textit{Validation:} This addresses the identified research gap that concept drift handling lacks automated strategy selection for quality signals \cite{contextdq2022}, grounded in the DTD theoretical framework \cite{dtd2026} proving that no single fixed adaptation strategy is universally optimal.

    \item \textbf{Online Autoencoder + Memory Module:} \textit{Problem:} Single-feature business rules cannot detect multivariate anomalies where individually reasonable feature combinations are collectively suspicious (e.g., a trip with valid fare and distance but impossible speed). \textit{Insight:} Denoising autoencoders learn a latent representation of normal data distributions; reconstruction error above a context-specific threshold indicates anomalous record structure. \textit{Method:} A denoising autoencoder trained on clean records learns a low-dimensional manifold of normal trip patterns. A FIFO memory module stores recent clean records, providing a reservoir for model warmup and reset. When ADWIN detects concept drift in quality meta-streams, the IEC triggers memory module reset: the autoencoder is warmup-retrained on the fresh clean records from the updated memory buffer, enabling the model to adapt to evolving distributions without catastrophic forgetting of the representation learned on historical data. This component detects the multivariate anomalies that single-feature business rules cannot express, complementing the 4D thresholding approach by operating on a different feature subspace. \textit{Validation:} Based on MemStream \cite{bhatia2022memstream} with extensions for zone-aware and periodicity-aware quality profiles, addressing the gap where memory-based methods lack geographic and temporal context for urban mobility data \cite{contextdq2022}.
\end{itemize}

\section{Thesis Content}
The remainder of this thesis is organized as follows. Chapter 1 provides background on streaming data processing (Kafka, Flink), data quality dimensions (DAMA-DMBOK), anomaly detection, concept drift, memory-augmented neural networks, and context-aware systems. Chapter 2 presents related work on streaming DQ monitoring, concept drift handling in ML systems, context-aware anomaly detection, and memory-based streaming methods, concluding with a qualitative comparison table that identifies the research gaps CA-DQStream addresses. Chapter 3 details the CA-DQStream system, including the approach direction, exploratory data analysis on the NYC Taxi dataset, and the proposed four-stage processing pipeline (Feature Engineering, Baseline Validation, Memory-Based Streaming Anomaly Detection, and Meta-Aggregation with the Intelligent Evolution Controller and Fault Tolerance mechanisms). Chapter 4 presents experimental evaluation organized by research question, validating all hypotheses on the NYC Yellow Taxi dataset with 72 million records spanning 26 months. Chapter 5 concludes with a summary of contributions, honest limitations, and future research directions.
